% Imporditud: tools/import_eksperimendid.py
% Allikas: Eesti füüsikaolümpiaadi eksperimendi ülesannete
%   kogu (Rael Kalda, 2023), ülesanne Ü26, lahendus L26.
\setAuthor{EFO žürii}
\setRound{piirkonnavoor}
\setYear{2017}
\setNumber{P E1}
\setDifficulty{1}
\setTopic{staatika}
\setVahendid{5-sendised euromündid (5 tükki), teadaoleva massiga joonlaud (mj = ...)}
\setPunktid{12}
\setAllikas{Eksperimendi kogu Ü26, lahendus lk 43}
\setLahendusseis{kasitsi}

\prob{Mündid}
Mündid laotakse silindriliselt üksteise otsa. Kuna mündi pinnal on muster, jäävad müntide vahele tühimikud. Mitu protsenti moodustavad tühimikud müntide silindri koguruumalast? Mündi materjali tihedus $\rho$ = 8,0 g/cm$^3$.

\hint

\solu
\textit{Lahendus on kogumikus lk 43. Siia ei ole seda veel ümber kirjutatud, sest originaalis on tabel või joonis, mis PDF-i tekstikihist loetavalt välja ei tule.}
\probend
